\section{Experimental Design}
\subsection{Vehicle fleet and task}
The benchmark contains 30 clients: ten nominal, ten payload, and ten tire-degraded. Payload clients increase mass and yaw inertia, while tire-degraded clients reduce front tire cornering stiffness. The true plant uses nonlinear lateral tire forces
\begin{equation}
 F_y=\mu F_z\tanh\left(\frac{C_\alpha\alpha}{\mu F_z}\right),
\end{equation}
whereas identification and controller synthesis use the linear two-state model. Every client sees the same two lane-change maneuvers, with three calibration episodes per maneuver. A sampled client is admitted only if a fixed baseline calibration is feasible on the standardized multi-maneuver bank. The primary output is yaw-rate RMSE; sideslip is bounded by $6^\circ$ and steering rate by $0.2618$ rad per simulation step.

\subsection{Landscape-similarity study}
To test whether dynamics heterogeneity translates into calibration heterogeneity, clients are evaluated on a common Sobol design $\{\xi^{(m)}\}_{m=1}^{M}$. Pairwise landscape similarity is
\begin{equation}
 \rho_{ij}^{J}=\operatorname{Spearman}\left(\{J_i(\xi^{(m)})\}_{m=1}^{M},\{J_j(\xi^{(m)})\}_{m=1}^{M}\right).
 \label{eq:landscape_similarity}
\end{equation}
The key comparison is $\rho_{ij}^{J}$ versus $\distdyn_{ij}$. \todo{Insert exact common Sobol-grid size $M$.}

\begin{figure*}[t]
\centering
\includegraphics[width=0.98\textwidth]{figures/fig2_dynamics_vs_landscape_similarity.pdf}
\caption{Plant heterogeneity does not translate into comparable controller-calibration heterogeneity. Despite substantial differences in identified dynamics, the LQI-calibration landscapes remain highly rank-correlated across the evaluated fleet.}
\label{fig:landscape}
\end{figure*}

\subsection{Pre-registered confirmatory campaign}
The final campaign uses 20 entirely new fleet seeds that were not used for development. For every seed, 30 mature clients generate independent BO histories of 25 evaluations. Six target problems are held out per seed, two per family, giving 120 confirmatory targets. Each target excludes its own mature history and has access to the remaining $29\times25=725$ historical records at time zero.

Each target starts with $n_{\rm init}=2$ local Sobol evaluations and has a maximum budget of 12 target-local experiments. The three methods use identical target initialization and evaluation randomness. The primary metric is $\Nfive$. Secondary metrics are early oracle-relative regret
\begin{equation}
 r_i^\star(n)=\frac{J_i^{\rm best}(n)-J_i^\star}{J_i^\star}
\end{equation}
and fixed-budget success probabilities. Fleet seed is the replication unit; target metrics are averaged within seed before paired bootstrap confidence intervals and Wilcoxon signed-rank tests are computed.

\begin{table}[t]
\centering
\caption{Final confirmatory protocol.}
\label{tab:protocol}
\scriptsize
\begin{tabular}{@{}ll@{}}
\toprule
Quantity & Setting\\
\midrule
Mature clients per seed & 30 (10 per family)\\
Confirmatory seeds & 20 untouched seeds\\
Held-out targets per seed & 6 (2 per family)\\
Mature history per client & 25 independent BO evaluations\\
Source records per target & 725\\
Target initialization & $n_{\rm init}=2$\\
Target-local budget & 12 evaluations\\
Methods & independent / global / similarity\\
Primary metric & $\Nfive$\\
\bottomrule
\end{tabular}
\end{table}

